\section{Purpose and relationship to the earlier releases}
\label{a2:purpose}

The original GitHub research manuscript separates mathematical legitimacy,
identification of an observation law and the predictive value of a fitted
procedure \citep{ProjectOriginal}. Its first addendum examines a separate
background-nuisance continuation whose null accounting remains incomplete
\citep{ProjectFirstAddendum}. This second companion adds a different empirical
axis: whether magnitude-class history or a more flexible magnitude-productivity
law improves the declared predictive comparisons. It reports the two completed
follow-ups, their numerical evidence and their stopping decisions, preserving
the earlier scientific records.

The distinction between releases is substantive. The manuscript's observed
reconstruction selected productivity ratios \(r=0.4,0.7,0.7\) at its three
origins. The new matched reference candidates use \(r=0.7\) throughout and a
different, common training-only magnitude-rate estimator
(Table~\ref{a2:recipe}). They are separately specified follow-up experiments,
not replacements for the manuscript's historical score table. Both use the
same hash-pinned GeoNet selection as our observed reconstruction. That archive
identity does not establish identity with the unavailable earlier analysis
discussed in the manuscript's catalogue-provenance section.

\input{recipe_table}

\subsection{External framework and attribution}

Self- and mutually exciting point processes provide a natural setting for
event-triggered clustering \citep{Hawkes1971}; earthquake triggering models
include the ETAS tradition \citep{Ogata1988}. Davis, Baeumer and Wang's
fractional earthquake model supplies context for Mittag-Leffler triggering
with magnitude productivity \citep{DavisBaeumerWang2024}. Their multivariate
fractional construction represents magnitude ranges as subprocesses and
permits history-dependent magnitude proportions
\citep{DavisBaeumerWang2026}. We inspected its model definitions in the
separately identified author preprint, Section~2, equations~(2)--(4)
\citep{DavisBaeumerWangPreprint2024}; those locators do not assert final-journal
pagination. For positive linear marked Hawkes models, partitioned multivariate
constructions provide an approximating parametric description under stated
assumptions \citep{DavisKresinBaeumerWang2026}. No stationarity or
identification theorem from that setting is transferred to our nonlinear
inhibitory comparator. These sources motivate the question and
model language; they do not supply the empirical findings reported here.

Our low-dimensional two-class shared-kernel family simplifies some class
and kernel structure while changing mark support and nuisance treatment.
The consulted MDFHP preprint \citep[Section~2]{DavisBaeumerWangPreprint2024}
uses truncated within-bin magnitude densities, including an upper magnitude
of 10. Our working exponential excess law is
unbounded. Consequently, our complete joint model is not a strict submodel
of that specification. The present comparisons neither reproduce nor refute
the authors' Japan and Middle America fitted results. Their results and
software must be credited in their own scope, and the locally derived
summation identity below is a result for the family actually implemented.
